% !TEX root = ../thesis.tex

% Some magic for unnumbered chapter correctly showing in ToC and headers
\chaptermark{Preface}
\phantomsection{}
\addcontentsline{toc}{chapter}{Preface}
\chapter*{Preface}\label{ch:preface}

I chose this topic because it offered an opportunity to work on a problem that combines artificial intelligence, remote sensing, and a real scientific application. My primary academic interest has long been connected with artificial intelligence and deep learning, and I wanted my bachelor thesis to become not only a formal study requirement, but also a meaningful first step towards more serious research work. The topic was especially suitable for this purpose because it was connected to an existing scientific background, real \gls{lidar} datasets, and a previously published study that provided an important methodological point of reference.

During the preparation of the thesis, I gradually realized that the central challenge was not limited to selecting and training a neural network model. A substantial part of the work was related to understanding the data and defining a reproducible experimental workflow that would be meaningful in an archaeological context. This changed my initial understanding of the task. At the beginning, I underestimated the amount of work required to formulate clear objectives, organize the methodology, and keep the experimental process consistent.

One of the most important challenges was the reconstruction of the experimental pipeline. In order to preserve consistency across more than eighty experiments, it became necessary to rethink the entire workflow. This decision was not immediate. Some earlier experiments had to be treated as exploratory rather than final, because later methodological corrections showed that they were not fully comparable with the final experimental protocol.

This thesis taught me that scientific work requires planning, patience, and attention to detail. It is not sufficient to focus only on model performance or individual technical improvements. The validity of the result depends on the entire chain of decisions, from the formulation of the problem to the preparation of data and the interpretation of results. For this reason, the work also became an important learning experience in how to approach research in a more systematic and critical way.